\chapter{Head Lice in Pregnancy}

\section*{Definition and Overview}
Head lice in pregnancy refers to an infestation of head lice that occurs while a woman is pregnant. Head lice are common at all ages, and pregnant women are just as susceptible, often acquiring them from children in the household. The condition and its treatment are essentially the same as at other times, but extra care is needed in choosing treatments, because some head lice products are not recommended during pregnancy. Most treatments, including thorough wet combing and certain lotions, are safe options for pregnant women. Head lice do not harm the baby and are not a sign of poor hygiene.

\section*{Epidemiology}
Head lice are extremely common, and women of childbearing age frequently encounter them, especially if they have children at primary school. It is estimated that around 1 in 10 children has head lice at any time, and household transmission to parents is common. There is little specific Australian data on head lice in pregnancy, but the condition is common enough that most antenatal care providers are asked about it at some stage. Infestation is not related to pregnancy itself.

\section*{Aetiology and Pathophysiology}
The biology of head lice is the same in pregnancy as at other times.

\begin{itemize}
    \item \textbf{The louse:} head lice are small insects that live on the scalp and feed on blood
    \item \textbf{Transmission:} by direct head-to-head contact, most often with children
    \item \textbf{Life cycle:} eggs (nits) hatch in 7--10 days, and adults lay new eggs
    \item \textbf{Itching:} an allergic reaction to louse saliva causes itching, which may take weeks to appear
    \item \textbf{No harm to pregnancy:} head lice stay on the scalp and do not affect the developing baby
    \item \textbf{No disease:} head lice do not transmit diseases
    \item \textbf{Treatment considerations:} the choice of treatment during pregnancy is guided by safety of the medicines
\end{itemize}

\section*{Clinical Features}
The features are the same as in any person with head lice.

\begin{itemize}
    \item Itching of the scalp, especially behind the ears and at the nape of the neck
    \item Live lice, which are small and move quickly
    \item Nits: eggs attached to the hair shaft near the scalp
    \item Scratch marks and sometimes secondary infection of the scalp
    \item Irritation or a scratchy sensation
    \item No symptoms in some people
\end{itemize}

\section*{Differential Diagnosis}
Other scalp conditions should be considered.

\begin{itemize}
    \item Dandruff, which flakes off easily (unlike nits, which are glued to the hair)
    \item Dry scalp and product residue
    \item Eczema and seborrhoeic dermatitis
    \item Insect bites
    \item Scabies, which affects other parts of the body
\end{itemize}

\section*{When to Seek Medical Care}
Pregnant women with head lice should seek advice from a pharmacist or doctor before choosing a treatment, to confirm the safest option in pregnancy. Medical care is also appropriate for secondary infection of the scalp.

\textbf{Call 000 (or local emergency services) for:}
\begin{itemize}
    \item This condition is not an emergency; however, seek care for spreading redness, swelling, or fever from an infected scalp
    \item Any severe allergic reaction to a treatment
    \item Obstetric emergencies are unrelated to head lice; seek urgent care for bleeding, pain, or reduced fetal movements as usual
\end{itemize}

\section*{Diagnosis and Investigations}
Diagnosis is made by finding live lice or nits.

\begin{itemize}
    \item \textbf{Wet combing:} combing conditioned hair with a fine-toothed comb reveals live lice
    \item \textbf{Nit examination:} nits are glued to the hair, distinguishing them from dandruff
    \item \textbf{Checking household contacts:} children and partners are checked and treated
    \item \textbf{Pharmacist or doctor advice:} to select a pregnancy-safe treatment
    \item \textbf{Follow-up:} repeat checks confirm clearance
\end{itemize}

\section*{Management}
Treatment during pregnancy emphasises safety and effectiveness.

\begin{itemize}
    \item \textbf{Wet combing:} a safe, chemical-free method using conditioner and a fine-toothed comb every 2--3 days for two weeks; this is often the preferred first option in pregnancy
    \item \textbf{Dimeticone lotion:} a physical (non-insecticidal) product that suffocates lice; it is generally considered safe in pregnancy, but confirm with a pharmacist
    \item \textbf{Permethrin:} a low-absorption insecticide that is considered safe in pregnancy when used as directed; check with a pharmacist or doctor
    \item \textbf{Avoiding some products:} certain insecticidal treatments are not recommended in pregnancy; a pharmacist can advise
    \item \textbf{Treating contacts:} children and partners are treated with age-appropriate products
    \item \textbf{Follow-up:} repeat combing and checks confirm the lice are gone
    \item \textbf{No fumigation needed:} extensive household cleaning is unnecessary
\end{itemize}

\section*{Complications}
\begin{itemize}
    \item Secondary infection of scratched skin
    \item Persistent infestation if treatment is not completed or contacts are not treated
    \item Anxiety about using treatments during pregnancy
    \item Skin irritation from treatments
    \item No complications for the pregnancy itself; head lice do not harm the baby
\end{itemize}

\section*{Prognosis and Outcomes}
Head lice in pregnancy resolve completely with appropriate treatment, and the outcome is excellent for both mother and baby. Wet combing and pregnancy-safe lotions clear the infestation within one to two weeks. Head lice do not affect the pregnancy, and the baby is not at risk. With proper treatment and checking of contacts, recurrence is prevented, and the condition causes no lasting harm.

\section*{Prevention and Self-Care}
\begin{itemize}
    \item Check your hair and your children's hair regularly
    \item Avoid sharing combs, brushes, hats, and hair accessories
    \item Treat promptly when live lice are found
    \item Ask a pharmacist or doctor which treatment is safest in pregnancy
    \item Consider wet combing as a chemical-free first option
    \item Check and treat all household members at the same time
    \item Inform the school or childcare centre
    \item Do not use products not recommended for pregnancy
    \item Attend routine antenatal care; head lice do not require special obstetric follow-up
\end{itemize}

\section*{Sources and Further Reading}
The following reputable sources provide further information for healthcare students and patients seeking reliable, current advice about head lice in pregnancy.

\begin{itemize}
    \item Pregnancy, Birth and Baby (Australian Government). \textit{Head lice during pregnancy.} https://www.pregnancybirthbaby.org.au/
    \item Healthdirect Australia. \textit{Head lice and safe treatment.} https://www.healthdirect.gov.au/
    \item NSW Health. \textit{Head lice fact sheet.} https://www.health.nsw.gov.au/
    \item Therapeutic Guidelines. \textit{Head lice treatment in pregnancy.} https://www.tg.org.au/
    \item MotherToBaby. \textit{Lice treatments and pregnancy.} https://mothertobaby.org/
    \item NHS. \textit{Head lice and nits.} https://www.nhs.uk/
\end{itemize}
